\scanpage{484}

\margin{so the \deleted{definition}\inserted{description}
of (33) amounts to the description of the \unsure{sole homomorphism|first homomorphism}
\diagram{\begin{tikzcd}[ampersand replacement=\&]
\mathcal M_{0,3}^{\sim}[0]\arrow[r]\&B'_0(\widehat{\mathbb Z})
\end{tikzcd}}}
Its \deleted{\unsure{not}} \unsure{\uline{value}} depends only on the pair
$(\alpha,\beta)$ \deleted{\unsure{part|taken}} of the triple $(\alpha,\beta,f)$
which describes an element of $\mathcal M_{0,3}^{\sim}$), and
the \deleted{\unsure{posi}} definition of $\underline u$ in terms of
$\alpha,\beta$ --- and even only in terms
of the image of $\xi=\alpha\rho(\alpha)^{-1}$ in
$\mathrm{Sl}(2,\widehat{\mathbb Z})$
\deleted{\unsure{modulo} $\unsure{\pm1}$, \unsure{or}}
\inserted{and even} only in
$\mathrm{Sl}(2,\widehat{\mathbb Z})'$ ! ---
\margin{For one \unsure{normalizes} by the condition
$\xi\rho(\xi)\rho^2(\xi)=\unsure{\pm1|+1}$ !}
\unsure{to|has} been \unsure{made|made} essentially in \S45
or at least in \S46 $\overline{\underline{\mathrm{III}}}$, --- we have
\[
\underline u:=\begin{pmatrix}1/D&B/D\\0&1\end{pmatrix}
\qquad\begin{gathered}
\text{where }\xi=\begin{pmatrix}A&B\\C&D\end{pmatrix}\\
\text{(in }\mathrm{Sl}(2,\widehat{\mathbb Z})\text{)}
\end{gathered}
\tag{41}
\]
or again, in terms of
\[
\underline\alpha=\begin{pmatrix}a&b\\c&d\end{pmatrix}
\qquad\text{with }ad-bc\unsure{=|\in}\inserted{\varepsilon'\in}\{\pm1\}
\]
\[
\left\{\begin{aligned}
\inserted{\varepsilon'}D&=(c^2+cd+d^2)\deleted{\unsure{/\varepsilon}}\\
\inserted{\varepsilon'}B&=ac+bd+bc
\end{aligned}\right.
\qquad\deleted{\text{with }\unsure{\varepsilon'}}
\tag{42}
\]

This does not depend on the foundations not laid for profinite groups ---
the ``calculations in $\xi,\eta$'' are in fact sufficient.
It is thanks to the \unsure{description} of this homomorphism
that one can define
\[
\begin{gathered}
\mathcal M_{0,3}^{\sim}(0)\subset\mathcal M_{0,3}^{\sim}[0]\\
\begin{array}{c}
\Vert\\
\{\underline{\underline u}\mid\underline u\in T_0(\widehat{\mathbb Z})
\quad\text{i.e. }B=0\}\\
\qquad\text{i.e. }ac+bd+bc=0
\end{array}
\end{gathered}
\tag{43}
\]
(normalization condition),

\mspage{483}

\scanpage{485}
\mspage{545}

and obtain canonical decompositions into
$\tfrac12$-\unsure{direct|direct} products
\[
\mathcal M_{0,3}^{\sim}[0]\simeq\mathcal M_{0,3}^{\sim}(0)\mathbin{.}L_{0\pi}
\tag{44}
\]
\[
\mathcal M_{0,3}^{\sim}\simeq\mathcal M_{0,3}^{\sim}(0)\mathbin{.}\mathcal G_{0,3}^{+\wedge}.
\tag{45}
\]

2\textsuperscript{o}) To describe (34), since it is already known on
$\mathcal G_{0,3}^{+\wedge}$, it suffices to describe it on
$\mathcal M_{0,3}^{\sim}[0]$ or even only on
$\mathcal M_{0,3}^{\sim}(0)$ (subject of course to compatibilities\ldots).
And on $\mathcal M_{0,3}^{\sim}[0]$ it coincides with (33) ---
whereas on $\mathcal M_{0,3}^{\sim}(0)$ it is practically trivial
and reduces to the multiplier $\mu$
\[
\Theta_G^{\pi}=\Theta_{B'_0}^{\pi}
\quad\text{on }\mathcal M_{\unsure{\rho|0},3}^{\sim}[0]
\tag{46}
\]
\[
\left\{\begin{aligned}
\Theta_G^{\pi}(\underline{\underline u})
  &=i'_{T_0}(\mu)
  \qquad\begin{gathered}\text{where }\mu\text{ is the}\\
  \text{multiplier of }\underline{\underline u}\end{gathered}\\
\text{\unsure{if|where}}\quad
\underline{\underline u}&\in\mathcal M_{0,3}^{\sim}(0)
\end{aligned}\right.
\tag{47}
\]
(this is even a possible definition of $\mathcal M_{0,3}^{\sim}(0)$).

\uline{Remark} As $\Theta_G^{\pi}$ was constructed
``by sheer strength of \unsure{wrist|wrists}'' by way
of $\Theta_{B'_0}^{\pi}$, i.e. $\underline U$ was constructed from $\underline u$,
there was for a long time a sort of

\mspage{484}

\scanpage{486}

confusion in my mind between $\underline u$ and $\underline U$,
all the more so since most of the time I restricted myself to working
in $\mathcal M_{\unsure{0|\rho},3}^{\sim}[0]$
(after having fumbled about quite a bit with more
\unsure{varied} points of view, before realizing that technically this was
the most elegant way to proceed --- from \S46 on).
This confusion was cleared up for good only yesterday and today,
by dint of putting my feet in my own dishes again and again\ldots

3\textsuperscript{o}) The homomorphisms (35)(36) were obtained
only by working in depth on the parameters $\alpha,\beta$
\unsure{for} $\underline u$ --- here again, it suffices to work
in $\mathcal M_{0,3}^{\sim}[0]$, and to define
\[
\diagram{\begin{tikzcd}[ampersand replacement=\&,row sep=tiny,column sep=huge]
\underline{\underline u}\arrow[r,mapsto]\&\underline a\\
\mathcal M_{0,3}^{\sim}[0]\arrow[r]\&T_{\rho}(\widehat{\mathbb Z})/\pm1\\
{}\arrow[r]\&\unsure{N'_{\sigma}}(\widehat{\mathbb Z})/\pm1\\
\underline{\underline u}\arrow[r,mapsto]\&\underline b
\end{tikzcd}}
\]

\mspage{485}

\scanpage{487}
\mspage{546}

so that it is trivial on $L_0^{\deleted{\unsure{\xi|\wedge}}\pi}$.

The quantities $\underline a,\underline b$ appeared in the
calculations gradually starting with
\S44 (page 366$'$ and following), light began to be shed on their
role only
starting with \S48 $\overline{\underline{\mathrm{III}}},
\overline{\underline{\mathrm{IV}}},\overline{\underline{\mathrm{V}}}$ ---
and a complete understanding of their nature, from the point
of view of the group schemes that enter into the theory,
likewise came about only in these very last
days. This understanding moreover remains for the moment rather
formal, and one should go deeper
in several directions,

a) Make explicit the meaning from the point
of view of the class fields of
$\mathbb Q(j)$ and \inserted{of} $\mathbb Q(\deleted{j}\inserted{i})$.

b) Make explicit the meaning for
the theory of ``\unsure{congruential|congruential}'' maps
(i.e. \unsure{subordinate|subordinate} to \deleted{\unsure{continuous}} representations
of the quotient $\mathrm{Sl}(2,\widehat{\mathbb Z})'$ of
$\mathrm{Sl}(2,\mathbb Z)^{\wedge\prime}$)
of the results obtained,

\mspage{486}

\scanpage{488}

c) Introduce on $\mathcal M_{\rho,\sigma}$ the elements of structure characteristic of arithmetic Galois groups -- and make the connection with the typical constructions of this sort in class field theory.

d) Examine to what extent in this context one can bring in the seemingly draconian compatibility conditions on $\alpha,\beta$ of §47 linked to compatibility with
\[
\diagram{\begin{tikzcd}[ampersand replacement=\&,column sep=large]
\varphi:\widehat\Pi'_{0,3} \arrow[r,"\sim"] \& \mathcal G_{0,3}^{+\wedge}
\end{tikzcd}}
\]
-- I \unsure{foresee.} moreover that the framework of the preceding calculations is too narrow to manage to bring these additional conditions into it.

e) What interests me most for the moment is a systematization of the preceding calculations in a framework much broader than that of the group $SG$ alone and of the \unsure{standard} homomorphism

\mspage{487}

\scanpage{489}
\mspage{547}

\[
\diagram{\begin{tikzcd}[ampersand replacement=\&,column sep=large]
\mathcal G_{0,3}^{+\wedge} \arrow[r] \& \deleted{\unsure{S}}G(\widehat{\mathbb Z})/\pm1
\end{tikzcd}}
\]
-- I have the feeling of having put my finger on a systematic procedure for constructing ``motivic'' representations of $\Gamma_{\mathbb Q}$, via representations of the ``geometric'' part $\mathcal G_{0,3}^{+\wedge}$ of
\[
\mathcal M_{0,3}^{+}\simeq\Pi_1(M_{0,3\,\mathbb Q},\unsure{\mathfrak S}_3)
\simeq\Pi_1(M'_{1,1\,\mathbb Q})
\]
-- or better, \inserted{via} representations of the discrete group $\mathcal G_{0,3}^{+}\simeq\mathrm{Sl}(2,\mathbb Z)'$ (or the \unsure{``}variants such as $\mathrm{Sl}(2,\mathbb Z)$, $\Pi_{0,3}$ \ldots) with values in the integral points of group schemes (\inserted{\unsure{generically}} reductive ??) over $\mathbb Z$.

f) Yet there is another direction of investigation that attracts me a great deal, not necessarily linked to the introduction of group schemes, and which could be pursued \inserted{\unsure{in vitro|in addition}} within the framework of just the ``big'' profinite fundamental groups --

\mspage{488}

\scanpage{490}

namely the study of certain special \inserted{$\mathbb Q$-}coverings of $\unsure{U}_{0,3\,\mathbb Q}=\mathbb P_{\unsure{0,3}}-\{0,1,\infty\}$, and information one should draw from them about additional properties of the ``arithmetic'' outer actions of $\Pi_{0,3}^{\wedge}$ (of the type of the results obtained in §47), and \deleted{\unsure{on}} calculations of splittings of the extension $\mathcal M_{0,3}$ of $\Gamma_{0,3}$ by $\mathcal G_{0,3}^{+\wedge}$ at points other than the ``critical'' points $0,1,\infty,\frac12,2,-1,-\bar j,\unsure{-j}$.

(Not counting the question \unsure{of establishing|of securing} the foundations, via the method suggested by Deligne -- and of studying the loop equation in the \unsure{protilp.|pro. tilp.} compactification $\unsure{L_A|L_\Delta}$ -- and of continuing to investigate in the direction of the \unsure{fundamental} conjectures, notably on the \unsure{arithmetization} of \uline{complex} points -- etc etc).

\mspage{489}

\scanpage{491}
\mspage{548}

Let us return for the moment to the question of identifying the \inserted{four} homomorphisms obtained (33) to (36), in terms of classical things. One can consider the homomorphism (33) \inserted{$\Theta_{B'_0}^{\pi}$} as being of an essentially technical nature, or at least as being subordinate to the homomorphism $\Theta_G^\pi$ of (34). Moreover, (35) and (36) are deduced from it, by composing \inserted{them} with the two homomorphisms
\[
\diagram{\begin{tikzcd}[ampersand replacement=\&,column sep=large]
\Gamma_{0,3}^{\sim} \arrow[r] \& \mathcal M_{0,3}^{\sim}
\end{tikzcd}}
\]
\inserted{(morally at least, for the moment !)} associated with the \deleted{\unsure{the}} points $\frac12$ and $-\bar j$ of $\unsure{U}_{0,3\,\mathbb Q}(\mathbb C)$, and which correspond, at the level of group schemes, to the two splittings $\psi_\rho,\psi_\sigma$ of $\overset{\inserted{(3)}}{\mathcal M_{\rho,\sigma}}$ over $\overset{\inserted{(31)}}{\mathcal T_{\rho,\sigma}}$, given \deleted{by} respectively by
\diagram{\begin{equation*}
\tag{48}
\left\{\begin{gathered}
\begin{tikzcd}[ampersand replacement=\&,row sep=small,column sep=large]
\psi_\rho:(\underline a,\underline b) \arrow[r,mapsto] \& (\underline a,\underline b,\underline a) \\
\psi_\sigma:(\underline a,\underline b) \arrow[r,mapsto] \& (\underline a,\underline b,\underline b)
\end{tikzcd}
\end{gathered}\right.
\end{equation*}}
\margin{\diagram{\begin{tikzcd}[ampersand replacement=\&,row sep=small,column sep=large]
\deleted{\unsure{\mathcal T_{\rho,\sigma}|\psi_\rho}}\,\mathcal T'_{\rho,\sigma}\arrow[r] \& \mathcal M_{\rho,\sigma} \\
\mathcal T_{\rho,\sigma}\arrow[r] \& \mathcal M_{\rho,\sigma},
\end{tikzcd}}}
with this grain of salt that the first is defined only on the subgroup $\mathcal T'_{\rho,\sigma}$ defined by
\[
\tag{49}
\left\{\begin{aligned}
\mathcal T'_{\rho,\sigma}&\subset\mathcal T_{\rho,\sigma}\\
\Vert\quad&\left\{\left(\underset{\inserted{\substack{\rotatebox[origin=c]{-90}{$\unsure{\in}$}\\T_\rho}}}{\underline a},\underset{\inserted{N'_\sigma}}{b}\right)\,\middle|\,\deleted{\text{in }\mathcal T_{\rho,\sigma}}\ \det\underline a=\det\underline b\right\}\\
&\Vert\quad T_\rho\mathbin{\underset{\mathbb{G}_m}{\times}}N'_\sigma.
\end{aligned}\right.
\]

\mspage{490}